
\medskip
\phantomsection\label{lem:sga3-VIB-8.1.2}%
\noindent\textbf{Lemma 8.1.2.---}
\emph{Let \(G\) be a sheaf of groups, and let \(A,B\) be two
group subpresheaves of \(G\).}
\begin{enumerate}
\item[\textup{\SGAInlineRefTarget{lem:sga3-VIB-8.1.2:item:i}(i)}]\phantomsection\label{concl:sga3-VIB-8.1.2-i}
  \emph{The morphism}
  \[
    \tau:\operatorname{Comm}(A,B)
      \longrightarrow
      \operatorname{Comm}(A^\flat,B^\flat)
  \]
  \emph{is a covering monomorphism.}

\item[\textup{(ii)}]\phantomsection\label{concl:sga3-VIB-8.1.2-ii}
  \emph{Consequently, there is an isomorphism}
  \phantomsection\label{form:sga3-VIB-8.1.2-CommT-isomorphism}%
  \[
    \operatorname{Comm}_T(A,B)
      \xrightarrow{\ \sim\ }
    \operatorname{Comm}_T(A^\flat,B^\flat).
  \]
\end{enumerate}

\emph{Proof.}
\textup{\SGARefLink{concl:sga3-VIB-8.1.2-i}{(i)}} Since \(A\) (respectively, \(B\)) is a
subpresheaf of \(G\), the associated sheaf \(A^\flat\)
(respectively, \(B^\flat\)) is a subpresheaf of \(G\)
containing \(A\) (respectively, \(B\)). It follows that
\(\tau\) is a monomorphism.

We now show that \(\tau\) is a covering. Let
\(S\in\operatorname{Ob}\mathcal C\) and let
\[
  g\in\operatorname{Comm}(A^\flat,B^\flat)(S).
\]
There exist an integer \(n\geq1\) and, for \(1\leq i\leq n\),
elements
\[
  a'_i\in A^\flat(S),\qquad
  b'_i\in B^\flat(S),\qquad
  \varepsilon_i\in\{\pm1\},
\]
such that
\[
  g=(a'_1,b'_1)^{\varepsilon_1}\cdots
    (a'_n,b'_n)^{\varepsilon_n},
\]
where \((a,b)\) denotes the commutator \(aba^{-1}b^{-1}\).
There is a refinement \(R\) of \(S\) such that
\[
  a'_i\in A(R)
  \qquad\text{and}\qquad
  b'_i\in B(R)
\]
for every \(i\). Then \(g\) is the composite morphism
\begin{center}
  \phantomsection\label{diag:sga3-VIB-8.1.2-composite}%
  \begin{tikzcd}[column sep=large]
    R
      \arrow[r,"{(a'_1,\ldots,b'_n)}"] &
    (A\times B)^n
      \arrow[r,"{\Phi^{\varepsilon_1,\ldots,\varepsilon_n}}"] &
    \operatorname{Comm}(A,B)
  \end{tikzcd}
\end{center}
where
\(\Phi=\Phi^{\varepsilon_1,\ldots,\varepsilon_n}\) is the
morphism defined set-theoretically by
\phantomsection\label{form:sga3-VIB-8.1.2-Phi}%
\[
  \Phi(a_1,b_1,\ldots,a_n,b_n)
    =(a_1,b_1)^{\varepsilon_1}\cdots
     (a_n,b_n)^{\varepsilon_n}
\]
for every \(T\in\operatorname{Ob}\mathcal C\), with
\(a_i\in A(T)\) and \(b_i\in B(T)\). This shows that
\[
  g\in\operatorname{Comm}(A,B)(R).
\]
It follows, as in the proof of IV, \SGARefLink{sga3:iv:proposition:4:3:11}{4.3.11}\textup{\SGARefLink{lem:sga3-VIB-8.1.2:item:i}{(i)}}, that
\[
  \operatorname{Comm}(A,B)
    \longrightarrow
  \operatorname{Comm}(A^\flat,B^\flat)
\]
is a covering.

The morphism
\[
  \operatorname{Comm}(A^\flat,B^\flat)
    \longrightarrow
  \operatorname{Comm}(A^\flat,B^\flat)^\flat
\]
is likewise a covering monomorphism
(IV, \SGARefLink{sga3:iv:proposition:4:3:11}{4.3.11}\textup{\SGARefLink{sga3:iv:proposition:4:3:11:item:iv:01}{(iv)}}). Hence the same is true of
\[
  \operatorname{Comm}(A,B)
    \longrightarrow
  \operatorname{Comm}(A^\flat,B^\flat)^\flat.
\]
By IV, \SGARefLink{sga3:iv:corollary:4:3:12}{4.3.12}, we therefore obtain an isomorphism
\[
  \operatorname{Comm}(A,B)^\flat
    \xrightarrow{\ \sim\ }
  \operatorname{Comm}(A^\flat,B^\flat)^\flat.
\]
This proves \hyperref[lem:sga3-VIB-8.1.2]{Lemma 8.1.2}.
